\documentclass{article}
\usepackage[utf8]{inputenc}
\usepackage[T1]{fontenc}
\usepackage{geometry}
\geometry{
    a4paper,
    left=2.5cm,
    right=2.5cm,
    top=2.5cm,
    bottom=2.5cm
}
\usepackage{graphicx}
\usepackage{float}
\usepackage{hyperref}
\usepackage[english]{babel}
\usepackage{amsmath}
\usepackage{amssymb}
\usepackage{array}
\usepackage{booktabs}
\usepackage{caption}
\usepackage{subcaption}
\usepackage{enumitem}
\usepackage{fancyhdr}
\usepackage{titlesec}
\usepackage{url}
\usepackage{textcomp}
\usepackage{xcolor}
\usepackage{listings}
\usepackage{multirow}
\usepackage{longtable}
\usepackage{tabularx}
\usepackage{times}

% Set up page style
\pagestyle{fancy}
\fancyhf{}
\rfoot{\thepage}
\renewcommand{\headrulewidth}{0pt}

% Title formatting
\titleformat{\chapter}[display]{\normalfont\bfseries}{\chaptertitle\thechapter}{0pt}{\Huge}
\titlespacing*{\chapter}{0pt}{50pt}{40pt}

% Document info
\title{Template Document}
\author{Academic Support Agent}
\date{\today}

\begin{document}
\maketitle

% First create the title page on its own.
\maketitlepage
\newpage

% Then create the table of contents on its own.

\newpage

% Finally create the abstract on its own.
\begin{abstract}
	\blindtext[3] % Replace this line with the abstract.
\end{abstract}
\newpage

% Paper content begins here.
\makeheader	% Adds the MLA-Style information
\maketitle	% Redefined to use MLA-Style

% Put paper here.
% Begin example text.
This is here just so I can cite this thingy \cite{Kavanaugh1976} and also this thingy \cite{Saussure1995} so there can be some examples of bibliography and such. Usually this is where you'd put your introduction or something.
\blinddocument
% End example text.

% Put the works-cited on a separate page.
\newpage
\makeworkscited
\end{document}